\documentclass[a4paper,twoside,11pt]{article}
\usepackage[utf8]{inputenc}
\usepackage[english]{babel}
\usepackage{graphicx}
\usepackage{url}

% pdflatex

% redefinição das margens das páginas
\setlength{\textheight}{24.00cm}
\setlength{\textwidth}{15.50cm}
\setlength{\topmargin}{0.35cm}
\setlength{\headheight}{0cm}
\setlength{\headsep}{0cm}
\setlength{\oddsidemargin}{0.25cm}
\setlength{\evensidemargin}{0.25cm}

\title{Prometheus Kotlin Client Library}

\author{
\begin{tabular}{c}
             Mário Rijo, n.º 50561, e-mail: A50561@alunos.isel.pt, tel.: 924422578\\
             Rafael Nicolau, n.º 50546, e-mail: A50546@alunos.isel.pt, tel.: 932456007\\
\end{tabular}}

\date{
\begin{tabular}{ll}
  {Supervisor:} & José Simão, e-mail: jose.simao@isel.pt \\
\end{tabular}\\
\vspace{5mm}
March 2025}

\begin{document}

\begin{figure}
\begin{center}
\resizebox{80mm}{!}{\includegraphics{logoISEL.png}}
\end{center}
\end{figure}

\maketitle

\section{Introduction}
In this project proposal there are some key contextual topics that are important to clarify.
The following paragraphs will explain the main context in detail.

\subsection{Prometheus}
Prometheus \cite{prometheus:prometheus} is an open-source metrics-focused monitoring and alerting toolkit originally built at SoundCloud in 2012. Prometheus has, since then, become one of the most popular monitoring tools, seeing major adoption across the industry and the open-source community.
It is now a standalone open-source project maintained independently of any company, joining the Cloud Native Computing Foundation in 2016 as the second hosted project after Kubernetes.

Prometheus is a pull-based monitoring system that scrapes metrics from instrumented jobs, either directly or via an intermediary push gateway for short-lived jobs.

\subsection{Prometheus Metric Types}
Prometheus provides four key metric types for system monitoring, each designed for specific performance tracking.

\subsubsection{Counter}
A Counter **only increases** or resets to zero on restart, useful for tracking requests or errors.

\subsubsection{Gauge}
A Gauge can **increase or decrease**, monitoring fluctuating values like memory usage or active connections.

\subsubsection{Histogram}
A Histogram groups values into **buckets** for analyzing distributions, useful for request and latency.

\subsubsection{Summary}
A Summary provides **precomputed quantiles** (e.g., median, 95th percentile) for precise performance tracking.

Each type enhances observability, allowing Prometheus to monitor and diagnose system performance effectively.

\subsection{Exporting a Library}
When dealing with observability tools like Prometheus, exporting plays a crucial role in making collected metrics accessible. A library designed for observability should provide mechanisms to expose data in a structured format that monitoring systems can scrape and analyze. In the case of Prometheus, this usually means:
\begin{itemize}
    \item Providing an HTTP endpoint where metrics can be queried.
    \item Formatting data according to the Prometheus exposition format.
    \item Supporting different metric types such as counters, gauges, and histograms.
\end{itemize}

For a Kotlin-based Prometheus client library, the export process generally involves defining metric types, registering them within the application, and exposing them through an endpoint that Prometheus can scrape at regular intervals. This allows real-time collection and monitoring of performance data, aiding in debugging and system optimization.

\section{The Project}

The goal of this project is to develop a Prometheus client library tailored for Kotlin, leveraging Kotlin's unique features such as coroutines. This library will provide seamless integration with Kotlin-based frameworks like Ktor, Http4k, and others, enabling efficient and scalable metric collection.

Our approach focuses on extending Prometheus functionality within the Kotlin ecosystem by designing idiomatic APIs that simplify metric instrumentation while maintaining high performance. By utilizing Kotlin’s lightweight concurrency model, the library will support asynchronous metric collection, making it well-suited for modern, event-driven applications.

Additionally, the library aims to provide extensions that facilitate effortless integration with popular Kotlin-based web frameworks, ensuring developers can easily expose and monitor application metrics. With a focus on usability, performance, and interoperability, this project will enhance observability in Kotlin applications while maintaining full compatibility with Prometheus’s existing ecosystem.

\section{Roadmap}
\begin{table}[h!]
\centering
\begin{tabular}{ |c||c|  }
  \hline
  Date &Task\\
  \hline
  March 10, 2025   & \textbf{Delivery - Proposal delivery}\\
  \hline
  March 17, 2025  & Finish studying KMP and the recommended structure for the library \\
  \hline
  April 14, 2025  &Implementation of at least two of the four metric types for all supported platforms\\
  \hline
  April 28, 2025&   \textbf{Presentation - Project progress}\\
  \hline
  May 12, 2025&   Implemention of the two remaining metric types for all supported platforms\\
  \hline
  May 26, 2025&   Integration with Popular Kotlin-based web frameworks\\
  \hline
  June 2, 2025 &\textbf{Delivery - Beta version}\\
  \hline
  June 23, 2025 &Testing and performance optimization\\
  \hline
  July 5, 2025 &Publication of the library on Maven Central\\
  \hline
  July 12, 2025    &\textbf{Delivery - Final version}\\
  \hline
 \end{tabular}
 \caption{Project's roadmap.}
 \label{table:2}
\end{table}

\bibliographystyle{plain}
\bibliography{refs}

\end{document}
